% GENERATED FILE. Edit canonical lessons, then run npm run generate:books.
% canonical-source-hash: fnv1a64:ef455e118f05b330
% canonical-lessons: MR-C334-kranti, MR-C334-sthapan-karne, MR-C334-ubharne, MR-C334-bakhar, MR-C334-silalekh, MR-R334-review-328-334

\chapter{Turning Points}
\label{ch:mr-turning-points}
\hlchaptermodality{334}

\begin{chapteropening}
\textbf{By the end of this chapter:} \emph{I can talk about great changes and what people set up.}

Complete the last lesson of chapter 334: I can talk about great changes and what people set up.
\end{chapteropening}

\section[\texorpdfstring{\mr{क्रांती} (krāntī)}{क्रांती (krāntī)}]{\mr{क्रांती} (krāntī) — a revolution}

\label{lesson:MR-C334-kranti}



\begin{quote}
Before the new one: say the Marathi for to set free, then the Marathi for to drive out.
\end{quote}

\subsection*{You'll want to know: \mr{क्रांती}}
\textbf{\mr{क्रांती}} — \emph{krāntī} — \textquotedblleft{}a revolution\textquotedblright{}.

\begin{quote}
\mr{नव्या} \mr{बियाण्यांमुळे} \mr{शेतीत} \mr{मोठी} \mr{क्रांती} \mr{झाली}.
\end{quote}

\begin{quote}
\emph{navyā biyāṇyāmmuḷe śetīt moṭhī krāntī jhālī.}
\end{quote}

\begin{quote}
New seeds brought about a revolution in farming.
\end{quote}

Feminine.

\mr{क्रांती} \mr{होणे}, a revolution happens.

\begin{cousinweb}
\textbf{\mr{क्रांती}} (\emph{krāntī}) is a tatsama, taken from Sanskrit \emph{krānti}, 'a step, a going; an overcoming' (Monier-Williams); the sense 'a revolution' is modern.
\end{cousinweb}

\begin{grammarlens}[title={What you've built}]
One more word for telling what happened.
\end{grammarlens}

\subsection*{Your turn}
\begin{itemize}
\raggedright
  \item \emph{Say it:} \emph{krāntī}
  \item \emph{Say it:} \emph{krāntī}, once more
  \item \emph{Say it:} \emph{navyā biyāṇyāmmuḷe śetīt moṭhī krāntī jhālī.} — New seeds brought about a revolution in farming.
  \item \emph{Recall:} say the Marathi for to set free, then the Marathi for to drive out, then say \emph{krāntī} again
\end{itemize}

\begin{tcolorbox}[breakable,colback=teal!4,colframe=teal!35!black,title={Before you move on}]
What does \mr{क्रांती} mean? (\textquotedblleft{}A revolution\textquotedblright{}.) Say it once more.
\end{tcolorbox}
\section[\texorpdfstring{\mr{स्थापन} \mr{करणे} (sthāpan karṇe)}{स्थापन करणे (sthāpan karṇe)}]{\mr{स्थापन} \mr{करणे} (sthāpan karṇe) — to found; to establish}

\label{lesson:MR-C334-sthapan-karne}



\begin{quote}
Before the new one: say the Marathi for to drive out, then the Marathi for a revolution.
\end{quote}

\subsection*{You'll want to know: \mr{स्थापन} \mr{करणे}}
\textbf{\mr{स्थापन} \mr{करणे}} — \emph{sthāpan karṇe} — \textquotedblleft{}to found; to establish\textquotedblright{}.

\begin{quote}
\mr{आजोबांनी} \mr{गावात} \mr{पहिली} \mr{शाळा} \mr{स्थापन} \mr{केली}.
\end{quote}

\begin{quote}
\emph{ājobānnī gāvāt pahilī śāḷā sthāpan kelī.}
\end{quote}

\begin{quote}
Grandfather founded the first school in the village.
\end{quote}

Past, with -\mr{ने} (he …): \mr{त्याने} \mr{स्थापन} \mr{केलं}, \emph{tyāne sthāpan kelaṁ}.

What is founded agrees: \mr{शाळा} \mr{स्थापन} \mr{केली}, \mr{संघ} \mr{स्थापन} \mr{केला}.

\begin{cousinweb}
\textbf{\mr{स्थापन} \mr{करणे}} (\emph{sthāpan karṇe}) joins \emph{sthāpan}, a tatsama taken from Sanskrit \emph{sthāpana}, 'a setting up, an establishing' (Monier-Williams), to \emph{karṇe}, 'to do'.
\end{cousinweb}

\begin{grammarlens}[title={What you've built}]
One more word for telling what happened.
\end{grammarlens}

\subsection*{Your turn}
\begin{itemize}
\raggedright
  \item \emph{Say it:} \emph{sthāpan karṇe}
  \item \emph{Say it:} \emph{sthāpan karṇe}, once more
  \item \emph{Say it:} \emph{ājobānnī gāvāt pahilī śāḷā sthāpan kelī.} — Grandfather founded the first school in the village.
  \item \emph{Recall:} say the Marathi for to drive out, then the Marathi for a revolution, then say \emph{sthāpan karṇe} again
\end{itemize}

\begin{tcolorbox}[breakable,colback=teal!4,colframe=teal!35!black,title={Before you move on}]
What does \mr{स्थापन} \mr{करणे} mean? (\textquotedblleft{}To found; to establish\textquotedblright{}.) Say it once more.
\end{tcolorbox}
\section[\texorpdfstring{\mr{उभारणे} (ubhārṇe)}{उभारणे (ubhārṇe)}]{\mr{उभारणे} (ubhārṇe) — to erect; to raise}

\label{lesson:MR-C334-ubharne}



\begin{quote}
Before the new one: say the Marathi for a revolution, then the Marathi for to found.
\end{quote}

\subsection*{You'll want to know: \mr{उभारणे}}
\textbf{\mr{उभारणे}} — \emph{ubhārṇe} — \textquotedblleft{}to erect; to raise\textquotedblright{}.

\begin{quote}
\mr{लोकांनी} \mr{चौकात} \mr{वीराचा} \mr{पुतळा} \mr{उभारला}.
\end{quote}

\begin{quote}
\emph{lokānnī caukāt vīrācā putḷā ubhārlā.}
\end{quote}

\begin{quote}
People put up the hero's statue in the square.
\end{quote}

Past, with -\mr{ने} (he …): \mr{त्याने} \mr{उभारलं}, \emph{tyāne ubhārlaṁ}.

What is raised agrees: \mr{पुतळा} \mr{उभारला}, \mr{इमारत} \mr{उभारली}.

\begin{cousinweb}
\textbf{\mr{उभारणे}} (\emph{ubhārṇe}) is a tadbhava: CDIAL lists Marathi \emph{ubhārṇe}, 'to raise', under Sanskrit \emph{udbharati}, 'raises' (CDIAL 2038).
\end{cousinweb}

\begin{grammarlens}[title={What you've built}]
One more word for telling what happened.
\end{grammarlens}

\subsection*{Your turn}
\begin{itemize}
\raggedright
  \item \emph{Say it:} \emph{ubhārṇe}
  \item \emph{Say it:} \emph{ubhārṇe}, once more
  \item \emph{Say it:} \emph{lokānnī caukāt vīrācā putḷā ubhārlā.} — People put up the hero's statue in the square.
  \item \emph{Recall:} say the Marathi for a revolution, then the Marathi for to found, then say \emph{ubhārṇe} again
\end{itemize}

\begin{tcolorbox}[breakable,colback=teal!4,colframe=teal!35!black,title={Before you move on}]
What does \mr{उभारणे} mean? (\textquotedblleft{}To erect; to raise\textquotedblright{}.) Say it once more.
\end{tcolorbox}
\section[\texorpdfstring{\mr{बखर} (bakhar)}{बखर (bakhar)}]{\mr{बखर} (bakhar) — a chronicle}

\label{lesson:MR-C334-bakhar}



\begin{quote}
Before the new one: say the Marathi for to found, then the Marathi for to erect.
\end{quote}

\subsection*{You'll want to know: \mr{बखर}}
\textbf{\mr{बखर}} — \emph{bakhar} — \textquotedblleft{}a chronicle\textquotedblright{}.

\begin{quote}
\mr{या} \mr{बखरीत} \mr{त्या} \mr{लढाईचं} \mr{वर्णन} \mr{आहे}.
\end{quote}

\begin{quote}
\emph{yā bakhrīt tyā laḍhāīcaṁ varṇan āhe.}
\end{quote}

\begin{quote}
This chronicle describes that battle.
\end{quote}

Feminine.

Old Marathi prose histories are called \mr{बखरी}.

\begin{cousinweb}
\textbf{\mr{बखर}} (\emph{bakhar}) is usually traced to Arabic \emph{khabar}, 'news, a report', through Persian; the Marathi form has moved its sounds round.
\end{cousinweb}

\begin{grammarlens}[title={What you've built}]
One more word for telling what happened.
\end{grammarlens}

\subsection*{Your turn}
\begin{itemize}
\raggedright
  \item \emph{Say it:} \emph{bakhar}
  \item \emph{Say it:} \emph{bakhar}, once more
  \item \emph{Say it:} \emph{yā bakhrīt tyā laḍhāīcaṁ varṇan āhe.} — This chronicle describes that battle.
  \item \emph{Recall:} say the Marathi for to found, then the Marathi for to erect, then say \emph{bakhar} again
\end{itemize}

\begin{tcolorbox}[breakable,colback=teal!4,colframe=teal!35!black,title={Before you move on}]
What does \mr{बखर} mean? (\textquotedblleft{}A chronicle\textquotedblright{}.) Say it once more.
\end{tcolorbox}
\section[\texorpdfstring{\mr{शिलालेख} (śilālekh)}{शिलालेख (śilālekh)}]{\mr{शिलालेख} (śilālekh) — an inscription}

\label{lesson:MR-C334-silalekh}



\begin{quote}
Before the new one: say the Marathi for to erect, then the Marathi for a chronicle.
\end{quote}

\subsection*{You'll want to know: \mr{शिलालेख}}
\textbf{\mr{शिलालेख}} — \emph{śilālekh} — \textquotedblleft{}an inscription\textquotedblright{}.

\begin{quote}
\mr{मंदिराच्या} \mr{भिंतीवर} \mr{एक} \mr{जुना} \mr{शिलालेख} \mr{आहे}.
\end{quote}

\begin{quote}
\emph{mandirācyā bhintīvar ek junā śilālekh āhe.}
\end{quote}

\begin{quote}
There is an old inscription on the temple wall.
\end{quote}

Masculine.

\mr{शिला}, a stone, + \mr{लेख}, a writing.

\begin{cousinweb}
\textbf{\mr{शिलालेख}} (\emph{śilālekh}) is built from Sanskrit \emph{śilā}, 'a stone', and \emph{lekha}, 'a writing' (Monier-Williams): writing in stone.
\end{cousinweb}

\begin{grammarlens}[title={What you've built}]
One more word for telling what happened.
\end{grammarlens}

\subsection*{Your turn}
\begin{itemize}
\raggedright
  \item \emph{Say it:} \emph{śilālekh}
  \item \emph{Say it:} \emph{śilālekh}, once more
  \item \emph{Say it:} \emph{mandirācyā bhintīvar ek junā śilālekh āhe.} — There is an old inscription on the temple wall.
  \item \emph{Recall:} say the Marathi for to erect, then the Marathi for a chronicle, then say \emph{śilālekh} again
\end{itemize}

\begin{tcolorbox}[breakable,colback=teal!4,colframe=teal!35!black,title={Before you move on}]
What does \mr{शिलालेख} mean? (\textquotedblleft{}An inscription\textquotedblright{}.) Say it once more.
\end{tcolorbox}
\section[(dialogue)]{Looking back: from War Comes to Turning Points}

\label{lesson:MR-R334-review-328-334}



\begin{quote}
Say the Marathi for an enemy and an attack.
\end{quote}

\subsection*{Your turn}
Say these aloud:

\begin{itemize}
\raggedright
  \item an enemy, an attack, an invasion, to fight, a battlefield
  \item a sword, a round shield, a weapon, a cannon, cavalry
  \item a commander, infantry, an army camp, to stand guard, to battle
  \item a siege, to surround, a moat, a rampart, a bastion
  \item a hero, a heroic deed, bravery, a sacrifice, to surrender
  \item an expedition, a military campaign, to retreat, to set free, to drive out
  \item a revolution, to found, to erect, a chronicle, an inscription
\end{itemize}

\begin{tcolorbox}[breakable,colback=teal!4,colframe=teal!35!black,title={Before you move on}]
An enemy? (\textbf{\mr{शत्रू}}, \emph{śatrū}.) An attack? (\textbf{\mr{हल्ला}}, \emph{hallā}.) An invasion? (\textbf{\mr{आक्रमण}}, \emph{ākramaṇ}.) To fight? (\textbf{\mr{लढणे}}, \emph{laḍhṇe}.) A battlefield? (\textbf{\mr{रणांगण}}, \emph{raṇāṅgaṇ}.) A sword? (\textbf{\mr{तलवार}}, \emph{talvār}.) A round shield? (\textbf{\mr{ढाल}}, \emph{ḍhāl}.) A weapon? (\textbf{\mr{शस्त्र}}, \emph{śastra}.) A cannon? (\textbf{\mr{तोफ}}, \emph{toph}.) Cavalry? (\textbf{\mr{घोडदळ}}, \emph{ghoḍdaḷ}.) A commander? (\textbf{\mr{सेनापती}}, \emph{senāpatī}.) Infantry? (\textbf{\mr{पायदळ}}, \emph{pāydaḷ}.) An army camp? (\textbf{\mr{छावणी}}, \emph{chāvṇī}.) To stand guard? (\textbf{\mr{पहारा} \mr{देणे}}, \emph{pahārā deṇe}.) To battle? (\textbf{\mr{झुंजणे}}, \emph{jhuñjṇe}.) A siege? (\textbf{\mr{वेढा}}, \emph{veḍhā}.) To surround? (\textbf{\mr{घेरणे}}, \emph{gherṇe}.) A moat? (\textbf{\mr{खंदक}}, \emph{khandak}.) A rampart? (\textbf{\mr{तटबंदी}}, \emph{taṭbandī}.) A bastion? (\textbf{\mr{बुरुज}}, \emph{buruj}.) A hero? (\textbf{\mr{वीर}}, \emph{vīr}.) A heroic deed? (\textbf{\mr{पराक्रम}}, \emph{parākram}.) Bravery? (\textbf{\mr{शौर्य}}, \emph{śaurya}.) A sacrifice? (\textbf{\mr{बलिदान}}, \emph{balidān}.) To surrender? (\textbf{\mr{शरण} \mr{जाणे}}, \emph{śaraṇ jāṇe}.) An expedition? (\textbf{\mr{मोहीम}}, \emph{mohīm}.) A military campaign? (\textbf{\mr{स्वारी}}, \emph{svārī}.) To retreat? (\textbf{\mr{माघार} \mr{घेणे}}, \emph{māghār gheṇe}.) To set free? (\textbf{\mr{मुक्त} \mr{करणे}}, \emph{mukta karṇe}.) To drive out? (\textbf{\mr{हाकलणे}}, \emph{hākalṇe}.) A revolution? (\textbf{\mr{क्रांती}}, \emph{krāntī}.) To found? (\textbf{\mr{स्थापन} \mr{करणे}}, \emph{sthāpan karṇe}.) To erect? (\textbf{\mr{उभारणे}}, \emph{ubhārṇe}.) A chronicle? (\textbf{\mr{बखर}}, \emph{bakhar}.) An inscription? (\textbf{\mr{शिलालेख}}, \emph{śilālekh}.)
\end{tcolorbox}
